% ============ 与 Deligne 的访谈 (原书 p.200, 2013 年 Abel 奖) ============
% 原书无编号,不出现在目录中; 此处用 \chapter* 保持该约定
\chapter*{与 Deligne 的访谈}

在宣布 2013 年 Abel 奖授予 Deligne 的仪式上，Gowers 给 Deligne 打了电话.下面是我对部分谈话的转写.

\textbf{Gowers:} 我的另一个问题.鉴于你对最后一个 Weil 猜想的证明显然绝对非凡，人们确实非常好奇，是什么让你觉得你有机会证明它.鉴于证明非常出人意料，很难理解你当时怎么知道它值得去做.

\textbf{Deligne:} 这是一个不错的故事.部分因为 Serre，也因为听了 Godement 的讲座，我对自守形式有些兴趣.Serre 明白 Ramanujan 猜想中的 $p^{11/2}$ 应该与 Weil 猜想本身有关系.Eichler 与 Shimura，以及 Verdier 已经做了很多工作，因此我理解了两者之间的联系.然后我读到 Rankin 的一些工作，他证明的不是想要的估计，而是差 $\frac{1}{4}$ 的东西---容易的结果与想要的相差 $\frac{1}{2}$.一旦我看到这样的东西，我就知道必须理解他在做什么，看看能否在其他情形做类似的事.于是我看了 Rankin，在那里我看到他用的是 Landau 的一个结果---思路是，当你有一个定义 $\zeta$ 函数的乘积时，你可以从 $\zeta$ 函数本身的极点信息得到局部因子的信息.极点由 Grothendieck 的理论在各种情形中相当容易给出.因此，用 Rankin 与 Landau 的这个方法，利用我们对极点的信息，看看能做什么，是很自然的.起初我不知道我能走多远.我能处理的第一个情形是射影空间中的奇维数超曲面.但那已经是一个全新的情形，于是我有信心可以一直走下去.这只是技巧问题.

\textbf{Gowers:} 听到这种事总是很好.当然，这传达了一个想法：最终通向这个惊人证明的是一系列自然的步骤.

\textbf{Deligne:} 是的，但要能看到这些步骤，关键是我不仅跟随代数几何的讲座，还跟随一些看起来相当不同的东西(现在会不那么不同)---自守形式理论.正是这两个领域能做到的事情之间的差距给出了必须做什么的答案.

\textbf{Gowers:} 那只是运气好，你恰好两样都知道吗？

\textbf{Deligne:} 是的.
